\section{Related Work}
\label{sec:related}

\subsection{Small Object Detection in High-Resolution Vision}
\label{ssec:related_sod}

Detecting tiny targets in high-resolution aerial and maritime imagery is severely challenged by extreme scale disparity, low contrast, and massive background areas ($>70\%\text{--}80\%$)~\cite{VisDrone,UAVDT,TinyPerson,SeaDronesSee}. Classical approaches utilize density-guided cropping (ClusDet~\cite{ClusDet}, DMNet~\cite{DMNet}) or scale-adaptive label assignment (NWD~\cite{NWD}, RFLA~\cite{RFLA}). Recently, ResBiDet~\cite{ResBiDet} introduced a coarse-to-fine framework using global heatmap guidance for region cropping. However, it relies on multi-scale image-level slicing and boundary stitching. In contrast, HESOD performs early feature-level selective routing directly within the detector trunk and explicitly optimizes instance coverage under irreversible pruning.

\subsection{Selective and Dynamic Computation}
\label{ssec:related_selective}

Dynamic computation accelerates inference by allocating compute strictly to prospective foreground regions (e.g., AutoFocus~\cite{AutoFocus}, QueryDet~\cite{QueryDet}, CEASC~\cite{CEASC}, and sparse location activation~\cite{Li2025Sparse}). Recently, coarse-to-fine detectors such as ESOD~\cite{ESOD} and OSDT~\cite{Wang2026OSDT} predict early spatial priority heatmaps to extract candidate patches. However, existing selectors rely exclusively on a single semantic score supervised by independent per-cell BCE. Under extreme scale disparity, micro targets suffer from semantic feature collapse after downsampling; meanwhile, per-cell BCE lacks joint instance-level coverage constraints, causing fatal false negatives under an asymmetric error cost. HESOD resolves this by combining semantic and spectral dual cues with object-level soft-coverage optimization.

\subsection{Spectral Saliency and Frequency Methods}
\label{ssec:related_spectral}

High-frequency edge gradients persist for micro targets even as photometric semantics fade~\cite{SET,Zhao2025TGRS}; recent methods exploit this via dense frequency transforms~\cite{SET}, frequency-gradient networks~\cite{Zhao2025TGRS}, dual-domain alignment~\cite{Guo2025DualDomain}, or multi-scale compensation~\cite{Wang2026HighFreq}, but all apply spectral cues \emph{downstream}, incurring substantial latency. In sharp contrast, HESOD repurposes high-frequency gradient saliency \emph{upstream} as an ultra-lightweight routing cue: channel pooling compresses stem features to 2 channels, enabling directional Laplacian and Sobel gradient filtering at constant width.


